%% kapa: exam=1 round=34 year=2023 period=1 form=A subject=civil_law subject_ko=민법 seq=25 number=25
%% source: Q-Net 감정평가사 기출문제 (cst00309) artlSeq=5211789 file=raw/1008/5211789/2023년도 제34회 감정평가사 1차시험 1교시.pdf
%% conversion: kordoc 4.24.51 -> tools/gen1.py (auto); review=proofread-sonnet; flags=none
\begin{question}{25}
동산의 선의취득에 관한 설명으로 옳지 않은 것은? (다툼이 있으면 판례에 따름)
\begin{choices}
  \item 등기나 등록에 의하여 공시되는 동산은 원칙적으로 선의취득의 대상이 될 수 없다.
  \item 선의취득이 성립하기 위해서는 양도인이 무권리자라고 하는 점을 제외하고는 아무런 흠이 없는 거래행위이어야 한다.
  \item 양도인이 제3자에 대한 반환청구권을 양수인에게 양도하고 지명채권 양도의 대항요건을 갖춘 경우, 선의취득에 필요한 점유의 취득 요건을 충족한다.
  \item 동산질권의 선의취득을 저지하기 위해서는 취득자의 점유취득이 과실에 의한 것임을 동산의 소유자가 증명하여야 한다.
  \item 양수인이 도품을 공개시장에서 선의·무과실로 매수한 경우, 피해자는 양수인이 지급한 대가를 변상하고 그 물건의 반환을 청구할 수 있다.
\end{choices}
\end{question}
\answer{4}
